% !TeX root = ../../Book3.tex
\section{Cauchy 积分公式}
\label{b3:sec:0203}

积分定理消去全纯被积函数的闭路积分。要从边界恢复内部一点 \(z\) 的函数值，就把被积函数改成 \(f(\zeta)/(\zeta-z)\)：它在 \(z\) 处留下一个不能消去的贡献。这个贡献可以压缩到任意小的圆周，再由连续性算出。本节先说明一般边界的消边步骤，随后研究对积分中的内部变量求导，以及由此产生的定量估计。

\subsection{基本公式}
\label{b3:subsec:020301}

前节的同伦证明适用于连续变形，不必预先给区域剖分。这里处理实际的区域边界时，使用多元微积分中的平面 Green 公式。所需版本如下：若 \(W\) 是有界平面区域，边界由有限条互不相交的分段光滑 Jordan 曲线组成，实函数 \(P,Q\) 在 \(\overline W\) 的开邻域为 \(C^1\)，则
\[
 \int_{\partial W}P\,\mathrm{d}x+Q\,\mathrm{d}y
 =\iint_W(Q_x-P_y)\,\mathrm{d}x\,\mathrm{d}y.
\]
这里 Jordan 曲线指除起终点相同外没有自交的闭曲线。边界取正向，即沿边界行走时区域在左侧；因此外边界逆时针，孔洞边界顺时针。第二卷定理 1.2.4 已给出 \(C^1\) 边界版本；本节采用上面允许分段光滑 Jordan 边界的经典 Green 公式作为实变量微积分先修，下面完整写出它在复积分中的应用。第一章的 Goursat 定理和全纯即解析并未依赖这个公式；此时全纯函数已经具有连续偏导数，才可使用这一版本。
% 跨卷背景定位：Book2/Part01/Chapter01/Section0102.tex，b2:thm:01-vector-green-plane，当前定理1.2.4；该处只承担C1边界版，不冒称已证明本节全部分段光滑Jordan边界版。

\Needspace{6\baselineskip}
\begin{lemma}[区域边界上的 Cauchy 定理]\label{b3:lem:02-cauchy-boundary}
设 \(W\subseteq\mathbb C\) 是有界区域，边界由有限条互不相交的分段光滑 Jordan 曲线组成，均按区域在左侧定向。若 \(h\) 在 \(\overline W\) 的一个开邻域上全纯，则
\[
 \int_{\partial W}h(\zeta)\,\mathrm{d}\zeta=0,
\]
其中左端表示沿所有边界分支的积分之和。
\end{lemma}
\begin{proof}
写 \(h=u+\mathrm{i}v\)、\(\zeta=x+\mathrm{i}y\)。第一章的解析性定理保证 \(u,v\) 为 \(C^1\)，而 Cauchy--Riemann 方程给出 \(u_x=v_y\)、\(u_y=-v_x\)。展开微分式并分别应用实 Green 公式，有
\[
 \begin{aligned}
 \int_{\partial W}h(\zeta)\,\mathrm{d}\zeta
 &=\int_{\partial W}(u\,\mathrm{d}x-v\,\mathrm{d}y)
   +\mathrm{i}\int_{\partial W}(v\,\mathrm{d}x+u\,\mathrm{d}y)\\
 &=\iint_W(-v_x-u_y)\,\mathrm{d}x\,\mathrm{d}y
   +\mathrm{i}\iint_W(u_x-v_y)\,\mathrm{d}x\,\mathrm{d}y=0.
 \end{aligned}
\]
两个面积积分的被积函数分别为零。孔洞的作用已经包含在正向边界的约定中，不需要把内边界遗漏或改成与外边界同向。
\end{proof}

这一引理允许直接在一个区域内挖去小圆盘。挖去的圆盘应当闭包含于原区域，函数则须在留下的闭区域附近全纯；原来的奇点可以落在挖去部分。

\Needspace{6\baselineskip}
\begin{theorem}[Cauchy 积分公式]\label{b3:thm:02-cauchy-formula}
设 \(V\subseteq\mathbb C\) 是有界单连通区域，其边界是一条正向分段光滑 Jordan 曲线，\(f\) 在 \(\overline V\) 的一个开邻域上全纯。则对每个 \(z\in V\)，
\[
 f(z)=\frac1{2\pi\mathrm{i}}\int_{\partial V}
       \frac{f(\zeta)}{\zeta-z}\,\mathrm{d}\zeta.
\]
\end{theorem}
\begin{proof}
固定 \(z\in V\)，选 \(\varepsilon>0\) 使 \(\overline{D(z,\varepsilon)}\subset V\)。函数 \(h(\zeta)=f(\zeta)/(\zeta-z)\) 在挖去这个圆盘后留下的闭区域附近全纯。前一引理给出
\[
 \int_{\partial V}\frac{f(\zeta)}{\zeta-z}\,\mathrm{d}\zeta
 =\int_{|\zeta-z|=\varepsilon}
       \frac{f(\zeta)}{\zeta-z}\,\mathrm{d}\zeta,
\]
右侧的小圆此时取逆时针方向。负号已经在将顺时针内边界移到等式右侧时消去。

小圆的参数式为 \(\zeta=z+\varepsilon e^{\mathrm{i}t}\)、\(0\le t\le2\pi\)，所以右侧等于
\[
 \mathrm{i}\int_0^{2\pi}f(z+\varepsilon e^{\mathrm{i}t})\,\mathrm{d}t.
\]
它与 \(2\pi\mathrm{i}f(z)\) 的差的模不超过
\[
 2\pi\max_{|\zeta-z|=\varepsilon}|f(\zeta)-f(z)|,
\]
由 \(f\) 在 \(z\) 连续而趋于零。左侧不依赖 \(\varepsilon\)，令 \(\varepsilon\downarrow0\) 即得公式。
\end{proof}

核的大小在小圆上为 \(1/\varepsilon\)，圆长却为 \(2\pi\varepsilon\)，两者抵消后留下有限的点值贡献。因此不能仅因小圆周长趋零就把其积分删去。若将核换成在 \(z\) 附近有界的函数，长度估计才会使这项积分趋于零。

\ProseParagraphBreak

“在闭区域的邻域全纯”保证边界积分有定义，也保证挖洞后可以应用 Green 公式。若只知道 \(f\) 在 \(V\) 内全纯，仍可在闭包含于 \(V\) 的小圆盘上使用第一章的公式；在没有边界值或边界长度不受控制时，不能直接把围道推到 \(\partial V\)。例如 \(1/(1-z)\) 在单位圆盘全纯，但单位圆上的被积式在 \(z=1\) 已经没有通常的连续曲线积分意义。

\ProseParagraphBreak

外边界公式也能计算不以所求点为中心的围道。例如在正向矩形
\(-2<\operatorname{Re}\zeta<2\)、\(-1<\operatorname{Im}\zeta<1\) 的边界上，
\[
 \int\frac{e^\zeta}{\zeta-\tfrac12}\,\mathrm{d}\zeta
 =2\pi\mathrm{i}e^{1/2}.
\]
这里无需分别计算四条边；被积式在矩形内只有分母指定的一个异常点。若把分母改成 \(\zeta-3\)，被积函数就在整个闭矩形附近全纯，边界积分为零。点是否位于围道内部是公式的一项条件，不能从分母的形式省略。

\Needspace{6\baselineskip}
\begin{proposition}[圆环上的 Cauchy 公式]\label{b3:prop:02-annulus-cauchy}
设 \(a\in\mathbb C\)、\(0<r<R<\infty\)，\(f\) 在闭圆环
\(\{\,\zeta\in\mathbb C:r\le|\zeta-a|\le R\,\}\) 的开邻域上全纯。若 \(r<|z-a|<R\)，则
\[
 f(z)=\frac1{2\pi\mathrm{i}}\int_{|\zeta-a|=R}
       \frac{f(\zeta)}{\zeta-z}\,\mathrm{d}\zeta
       -\frac1{2\pi\mathrm{i}}\int_{|\zeta-a|=r}
       \frac{f(\zeta)}{\zeta-z}\,\mathrm{d}\zeta,
\]
其中两个圆周都取逆时针方向。
\end{proposition}
\begin{proof}
在圆环内部挖去闭包含于其中的 \(D(z,\varepsilon)\)。对剩余区域使用\cref{b3:lem:02-cauchy-boundary}，外圆的积分减去内圆的积分，等于绕 \(z\) 的逆时针小圆积分。后一积分按前一定理的参数计算趋于 \(2\pi\mathrm{i}f(z)\)，即得结论。
\end{proof}

取 \(a=0\)、\(f(\zeta)=1/\zeta\)，可以把孔洞项算得很清楚。由于 \(z\ne0\)，
\[
 \frac1{\zeta(\zeta-z)}
 =\frac1z\left(\frac1{\zeta-z}-\frac1\zeta\right).
\]
在外圆上，括号中两项的积分都是 \(2\pi\mathrm{i}\)，故整个外圆积分为零。在内圆上，第一项在圆盘附近全纯，积分为零；第二项的积分为 \(2\pi\mathrm{i}\)，故整个内圆积分为 \(-2\pi\mathrm{i}/z\)。公式恰好由减去内圆贡献恢复 \(1/z\)。若只保留外圆，将错误地得到零。

\ProseParagraphBreak

圆环公式将在下一章产生 Laurent 展开：外圆适合按 \(z-a\) 的正幂展开核，内圆适合按它的负幂展开核。这里先把两条边界及其符号固定下来，后面的级数计算才有正确起点。

\subsection{导数公式}
\label{b3:subsec:020302}

积分公式把内部变量放在一个简单的有理核中。只要内部点与边界保持正距离，核的差商误差便可在整条边界上一致控制；这正是将求导移入积分号的依据。

\Needspace{6\baselineskip}
\begin{proposition}[Cauchy 导数公式]\label{b3:prop:02-cauchy-derivatives}
设 \(V\subseteq\mathbb C\) 是有界单连通区域，其边界是一条正向分段光滑 Jordan 曲线，\(f\) 在 \(\overline V\) 的开邻域上全纯。则对每个整数 \(n\ge0\) 和 \(z\in V\)，
\[
 f^{(n)}(z)=\frac{n!}{2\pi\mathrm{i}}\int_{\partial V}
        \frac{f(\zeta)}{(\zeta-z)^{n+1}}\,\mathrm{d}\zeta.
\]
\end{proposition}
\begin{proof}
固定非空紧集 \(K\subset V\)，记 \(d=\operatorname{dist}(K,\partial V)>0\)。当 \(z\in K\)、\(|h|<d/2\) 时，
\[
 \frac1h\left(\frac1{\zeta-z-h}-\frac1{\zeta-z}\right)
 =\frac1{(\zeta-z-h)(\zeta-z)}.
\]
右端与 \((\zeta-z)^{-2}\) 的差为
\[
 \frac{h}{(\zeta-z-h)(\zeta-z)^2},
\]
其模至多为 \(2|h|/d^3\)，对所有 \(\zeta\in\partial V\)、\(z\in K\) 同时成立。因此将差商移入积分后，误差的模不超过
\[
 \frac{L(\partial V)\cdot\max_{\partial V}|f|}{2\pi}
       \frac{2|h|}{d^3}\longrightarrow0.
\]
这证明一阶公式。

为使归纳步骤同样明确，令 \(m\ge1\)、\(w=\zeta-z\)。把实变量的二阶积分余项公式用于 \(t\mapsto(w-th)^{-m}\)，得到
\[
 (w-h)^{-m}-w^{-m}-mh\,w^{-m-1}
 =m(m+1)h^2\int_0^1(1-t)(w-th)^{-m-2}\,\mathrm{d}t.
\]
由 \(|w-th|\ge d/2\)，右端的模至多为
\[
 m(m+1)2^{m+1}\frac{|h|^2}{d^{m+2}}.
\]
除以 \(h\) 后误差一致趋零，所以核的导数可以再次移入积分。对 \(m=1,2,\ldots\) 依次进行，每步产生正因子 \(m\)，第 \(n\) 阶的总因子为 \(n!\)。任一内部点都可放入这样的紧集，结论因而成立于整个 \(V\)。
\end{proof}

求导的正号来自内部变量 \(z\)：对 \((\zeta-z)^{-1}\) 求导时，幂函数的负号与内层 \(-z\) 的负号相消。例如对包围 \(a\) 的正向圆周，只要 \(f\) 在闭盘附近全纯，
\[
 \int\frac{f(\zeta)}{(\zeta-a)^3}\,\mathrm{d}\zeta
 =\pi\mathrm{i}f''(a).
\]
取 \(f(\zeta)=\zeta e^\zeta\)，有 \(f''(a)=(a+2)e^a\)，于是积分等于 \(\pi\mathrm{i}(a+2)e^a\)。分母的三次幂对应二阶导数及因子 \(2!\)，不是三阶导数。

\ProseParagraphBreak

同一估计还说明积分可以用来构造全纯函数。设 \(\gamma:[\alpha,\beta]\to\mathbb C\) 为分段光滑曲线，\(g\) 在其像上连续。在开集 \(W=\mathbb C\setminus\gamma([\alpha,\beta])\) 上定义
\[
 H(z)=\int_\gamma\frac{g(\zeta)}{\zeta-z}\,\mathrm{d}\zeta.
\]
对每个非空紧集 \(K\subset W\)，距离 \(d=\operatorname{dist}(K,\gamma([\alpha,\beta]))\) 为正，刚才的差商估计仍然适用，故
\[
 H^{(n)}(z)=n!\int_\gamma
          \frac{g(\zeta)}{(\zeta-z)^{n+1}}\,\mathrm{d}\zeta.
\]
这里 \(g\) 只需沿曲线连续；所求导的变量出现在核中，并没有要求 \(g\) 延拓为全纯函数。

\ProseParagraphBreak

对正向单位圆与 \(g\equiv1\)，\(H\) 在圆内恒为 \(2\pi\mathrm{i}\)，在圆外恒为零。圆内的值来自积分公式；圆外的值来自被积函数在闭单位圆盘附近全纯。两个连通分支上的全纯性并不提供跨越单位圆的连续性。若在曲线上把数据 \(g\) 改为 \(\widetilde g\)，相应函数之差却有直接控制：
\[
 \sup_{z\in K}|H^{(n)}(z)-\widetilde H^{(n)}(z)|
 \le \frac{n!L(\gamma)}{d^{n+1}}
       \max_{\zeta\in\gamma}|g(\zeta)-\widetilde g(\zeta)|.
\]
因此边界数据的小误差能控制远离曲线处的各阶导数；离曲线越近或求导阶数越高，所需的误差控制越强。

\subsection{Cauchy 估计}
\label{b3:subsec:020303}

\term[cauchy-guji]{Cauchy 估计}[Cauchy estimates]把边界上的函数大小变成内部的导数大小。圆周特别适合这一步，因为中心到各边界点的距离相同。

\Needspace{6\baselineskip}
\begin{theorem}[Cauchy 估计]\label{b3:thm:02-cauchy-estimate}
设 \(a\in\mathbb C\)、\(r>0\)，\(f\) 在 \(\overline{D(a,r)}\) 的开邻域上全纯，记 \(M_r=\max_{|z-a|=r}|f(z)|\)。则对每个整数 \(n\ge0\)，
\[
 |f^{(n)}(a)|\le\frac{n!M_r}{r^n}.
\]
\end{theorem}
\begin{proof}
在中心的导数公式中，分母的模为 \(r^{n+1}\)，圆长为 \(2\pi r\)。长度估计给出
\[
 |f^{(n)}(a)|
 \le\frac{n!}{2\pi}\cdot2\pi r\cdot\frac{M_r}{r^{n+1}}
 =\frac{n!M_r}{r^n}.
\]
\end{proof}

若 \(f\) 只在 \(D(a,R)\) 上全纯且 \(|f|\le M\)，对任意 \(r<R\) 应用估计，再令 \(r\uparrow R\)，仍得
\[
 |f^{(n)}(a)|\le\frac{n!M}{R^n}.
\]
这个极限只涉及已经有共同上界的数值，不需要先定义最外圆周上的函数值。估计中的常数不能普遍减小：\(f(z)=M((z-a)/r)^n\) 在取样圆上模恒为 \(M\)，中心的 \(n\) 阶导数恰好达到上界。

\ProseParagraphBreak

把圆心移到待估计的点，就能得到紧集上的统一控制。具体地，设 \(K\) 为非空紧集，\(\delta>0\)，而 \(f\) 在
\[
 K_\delta=\{\,z\in\mathbb C:\operatorname{dist}(z,K)<\delta\,\}
\]
全纯且 \(|f|\le M\)。每个 \(a\in K\) 都有 \(D(a,\delta)\subset K_\delta\)，故 \(|f'(a)|\le M/\delta\)。还要估计邻近两点的函数差时，应确保连接它们的线段也处于可控范围。若 \(a,b\in K\)、\(|a-b|<\delta/2\)，线段上的每个点到 \(K\) 的距离小于 \(\delta/2\)，以它为中心的半径 \(\delta/2\) 圆盘包含于 \(K_\delta\)。于是沿线段有 \(|f'|\le2M/\delta\)，由端点公式，
\[
 |f(b)-f(a)|\le\frac{2M}{\delta}|b-a|.
\]
这说明一族函数若在同一邻域有共同上界，就在 \(K\) 上具有共同的连续性控制。第四章的紧性论证将使用这一机制。

\ProseParagraphBreak

半径有时也需要选择。设整函数满足 \(|f(z)|\le A e^{B|z|}\)，其中 \(A,B>0\)，写 \(f(z)=\sum_{n\ge0}c_nz^n\)。对每个 \(r>0\)，
\[
 |c_n|\le A\frac{e^{Br}}{r^n}.
\]
固定整数 \(n\ge1\)，右端的对数除常数外为 \(Br-n\log r\)。它的导数是 \(B-n/r\)，在 \(r<n/B\) 时为负，在 \(r>n/B\) 时为正，故最佳选择为 \(r=n/B\)，得到
\[
 |c_n|\le A\left(\frac{eB}{n}\right)^n.
\]
这里先固定求导阶数，再为这一阶选择圆周。对 \(f=e^z\) 得到 \(1/n!\le(e/n)^n\)：估计未恢复精确系数，却已经显示高阶系数快速衰减。后面研究整函数增长时，同样的优化会反复出现。

\Needspace{6\baselineskip}
\begin{proposition}[Taylor 截断误差]\label{b3:prop:02-taylor-remainder}
设 \(a\in\mathbb C\)、\(R>0\)、\(M\ge0\)，\(f\) 在闭圆盘 \(\overline{D(a,R)}\) 的邻域全纯，且边界上 \(|f|\le M\)。对整数 \(N\ge0\)，令 \(T_N(z)=\sum_{n=0}^N f^{(n)}(a)(z-a)^n/n!\)。若 \(0\le r<R\)，则
\[
 \sup_{|z-a|\le r}|f(z)-T_N(z)|
 \le \frac{M(r/R)^{N+1}}{1-r/R}.
\]
\end{proposition}
\begin{proof}
第一章保证 \(f\) 在圆盘内等于中心 Taylor 级数。Cauchy 估计给出 \(|f^{(n)}(a)|/n!\le M/R^n\)，故截断之后的尾项模不超过 \(M\sum_{n>N}(r/R)^n\)。求几何级数的和即得结论；\(r=0\) 时余项本身为零。
\end{proof}

这个界与第一章的误差估计一致，现在可以用来反推所需的截断次数。设 \(M>0\)、\(0<r<R\)，记 \(q=r/R\)，希望误差不超过 \(\varepsilon>0\)。充分条件是
\[
 N+1\ge
 \frac{\log\bigl(M/(\varepsilon(1-q))\bigr)}{\log(1/q)}.
\]
若右端为负，任意非负截断次数都满足该充分条件；否则选择满足不等式的整数即可。例如 \(r=R/2\)、\(M=1\)，余项界为 \(2^{-N}\)，取 \(N=10\) 就使误差小于 \(10^{-3}\)。若 \(r\) 接近 \(R\)，相同次数的保证会明显变差，因为几何衰减的比值接近一。

\ProseParagraphBreak

取 \(N=0\)，还得到 \(|f(z)-f(a)|\le M|z-a|/(R-|z-a|)\)。这与先估计导数再沿线段积分是同一局部控制的两种用法。无论采用哪一种，都要明确函数上界在哪个圆盘成立，以及待估计点与取样边界之间还留下多少距离。
